% Volume II, appendix B: the tools of the engine (integrator, 2026-10-16).
% The table of options is checked by memoire/outils/tome2_annexes.py: every option that
% moteur/bins/troncons/src/main.rs reads must appear in it, written \option{name}, in English and in
% French; an option too many or too few makes the tool fall, hence the build.
% The three other tables are generated by the same tool from git ls-files.
\newcommand{\option}[1]{\texttt{-{}-#1}}
\chapter{Tools of the engine}\label{t2ann:outils}

\begin{origine}
\textbf{Where we start.} Appendix B of Volume~I (\enquote{Tools}): every tool of the studies, sorted by study, with the first sentence of its own description, and the chain that builds the thesis.
\textbf{What this appendix establishes.} The same for the system: the chain that builds this volume and checks its numbers; the \texttt{troncons} program, its four subcommands and every one of its options, kept up to date by a witness; every committed tool of the five lots, of the trial and of the round; every figure and the files it reads; every folder of logs.
\textbf{What remains.} The trial options (\option{kv-f32}, \option{longueur-lecture}) await the arbiter's decision to become defaults or disappear.
\end{origine}

\section{The chain of Volume~II}\label{t2ann:sec:chaine}

Volume~II is built like the first, by the same script and the same \texttt{latexmk} configuration (\fichier{memoire/construire.sh}, \fichier{memoire/latexmkrc}):
\begin{itemize}
\item \texttt{./construire.sh tome2} (or \texttt{tome2-en}, \texttt{tome2-fr}) regenerates the data of Volume~II's figures, redoes its numbers by a second path, regenerates the tables of this appendix, then builds the two PDFs with \texttt{latexmk} (lualatex, biber) in \fichier{memoire/tome2/en/} and \fichier{memoire/tome2/fr/}; \texttt{./construire.sh} without argument builds both volumes. The script fails at the first error, at any missing reference or citation, at any entry that biber does not find or finds twice, and at any data tool that exits with an error.
\item \fichier{memoire/outils/tome2_format.py}: the CSV files of the \enquote{Architecture} and \enquote{Format} chapters (lines of code and witnesses per member, rates per tensor, chunks per layer, memory peaks of the coder), drawn from the \texttt{Cargo.toml} files and lot F's logs, and the calculations of the format chapter.
\item \fichier{memoire/outils/tome2_chemin_noyaux.py}: the CSV files of the chapters \enquote{The path of the bytes} and \enquote{The kernels}, drawn from the logs of lots E, K and O.
\item \fichier{memoire/outils/tome2_reference_cache.py}: the CSV files of the \enquote{Reference} and \enquote{Cache} chapters, and 31 published numbers redone from the logs; it stops at the first gap, and \texttt{--casser} must make it fall.
\item \fichier{memoire/outils/tome2_resultats.py}: the CSV files of the \enquote{Results}, \enquote{Trial} and \enquote{Discussion} chapters; it rereads every pass, recomputes means, noises, ratios, Wilson intervals and floors, and exits with an error if a recomputation departs from a published number.
\item \fichier{memoire/outils/tome2_annexes.py}: the tables of this appendix and the threshold table of \cref{t2ann:protocoles}, from \texttt{git ls-files} and \texttt{git log}; witness: the options of \cref{t2ann:tab:options} against those the code reads.
\end{itemize}

\section{The \texttt{troncons} program}\label{t2ann:sec:programme}

The program (\fichier{moteur/bins/troncons/src/main.rs}, lot O) has four subcommands. \texttt{generer} prefills the prompt, then produces $n$ greedy tokens, and writes the text and a JSON log (command, date, power, engine counters, tokens/s, peak memory and VRAM, bytes read from disk by the system counter and by the engine, fingerprint of the text). \texttt{mesurer} is \texttt{generer} with ten seconds of rest before the pass, to quantify the idle disk. \texttt{evaluer} runs the windows of a llama-perplexity \texttt{.kld} base file (same tokens, each from position 0), passes the scored logits through the same 16-bit quantisation as llama.cpp (criterion C4) and writes the fingerprint64 of the logits of every step (criterion C3). \texttt{disque} measures the disk stage alone, in the same session as the engine: $n$ chunks drawn without replacement with a published seed, read unbuffered at the engine's length, with \texttt{profondeur} reads in flight. \Cref{t2ann:tab:options} gives every option.

\begingroup\small
\renewcommand{\arraystretch}{1.1}
\begin{longtable}{@{}p{0.30\linewidth}p{0.17\linewidth}p{0.47\linewidth}@{}}
\caption[The options of the troncons program]{The options of the \texttt{troncons} program, as \fichier{moteur/bins/troncons/src/main.rs} reads them (checked by \fichier{memoire/outils/tome2_annexes.py}: none too many, none too few). The names are French, as in the code.}\label{t2ann:tab:options}\\
\toprule
Option & Default & Effect \\
\midrule
\endfirsthead
\toprule
Option & Default & Effect \\
\midrule
\endhead
\multicolumn{3}{@{}l}{\textsf{\bfseries The engine (every subcommand that opens it)}} \\
\option{modele} \emph{f.trc} & required & the \texttt{.trc} file served \\
\option{ram-tiede} \emph{GB} & 0 (mode A0) & warm cache: coded chunks kept in pinned memory \\
\option{enchassement-en-ram} & no & keep the embedding table (0.6\,GB) in pinned memory instead of reading one row per token \\
\option{profondeur} \emph{n} & 4 & reads entrusted to the system at once \\
\option{couches-d-avance} \emph{d} & 3 & layers whose arrival buffers rotate (window of $d \cdot k$ buffers) \\
\option{contexte} \emph{n} & \num{4096} & positions at most (KV cache of the full layers) \\
\option{temoin} & no & every decoded expert goes back to the host, its fingerprint compared to the directory (C2) \\
\option{brut} & no & the expected file is raw (flag bit 2); it must be \\
\option{emplacements} \emph{n} & from the VRAM & expert slots in video memory \\
\option{vram-plafond} \emph{MiB} & \num{7179} & ceiling of the total VRAM used to compute the slots: $8{,}188 - 909 - 100$ \\
\option{longueur-lecture} \emph{B} & the format's & length of every chunk read, a multiple of \num{4096} (trial: \cref{sec:t2c-regimes}) \\
\option{chronogramme} \emph{f.csv} & none & dates of every stage of every miss \\
\option{prechargement} & no & J1+: anticipated reads of the experts predicted by the next layer's router (tried, abandoned) \\
\option{sans-battement} & heartbeat & removes the line written at every layer of the streamed prefill (finding G1) \\
\option{sans-garde} & guard & removes the production guard; serves only to quantify its cost \\
\option{casser-garde} \emph{troncon:N} | \emph{ligne:N} & none & witness of the guard: flips one bit of the $N$-th chunk or $N$-th row read \\
\option{kv-f32} & f16 & trial: f32 KV cache (finding M5; seam N with the arbiter) \\
\multicolumn{3}{@{}l}{\textsf{\bfseries \texttt{generer}, \texttt{mesurer}}} \\
\option{vocab} \emph{f.vocab.gguf} & required & the tokenizer companion \\
\option{invite} \emph{f.txt} & required without \option{harmonie} & the prompt, in UTF-8 (its SHA-256 goes to the log) \\
\option{harmonie} \emph{\enquote{question}} & none & a prompt in gpt-oss's harmony conversation format \\
\option{systeme} \emph{text} & fixed message & the system message of the harmony prompt \\
\option{n} \emph{n} & 256 & tokens produced (\texttt{disque}: chunks read, \num{1000}) \\
\option{repos} \emph{s} & 0; 10 (\texttt{mesurer}) & seconds of rest before the pass, to quantify the idle disk \\
\option{dll} \emph{f.dll} & llama.cpp b11515 & the tokenizer library \\
\option{arret-fin} & no & stop at the end token of the harmony conversation \\
\option{sans-vram} & VRAM sampled & do not sample the VRAM \\
\option{flux} & step loop & streamed prefill (J1+): every expert read once for all positions \\
\option{journal} \emph{f.json} & none & the log of the pass \\
\option{sortie} \emph{f} & none; required by \texttt{evaluer} & the text produced; for \texttt{evaluer}, the \texttt{.kld} written \\
\option{trace} \emph{f.trace} & none & the cache trace (for \texttt{rejouer}) \\
\multicolumn{3}{@{}l}{\textsf{\bfseries \texttt{evaluer}}} \\
\option{base} \emph{f.kld} & required & llama-perplexity's base file (tokens and logits) \\
\option{fenetres} \emph{n} & all & windows evaluated \\
\option{positions} \emph{n} & all (512) & positions evaluated per window \\
\option{empreintes} \emph{f.csv} & none & the fingerprint64 of the f32 logits of every step (C3) \\
\multicolumn{3}{@{}l}{\textsf{\bfseries \texttt{disque}}} \\
\option{graine} \emph{g} & 1 & seed of the draw of chunks (xorshift64*, that of lot E's experiments) \\
\bottomrule
\end{longtable}
\endgroup

\section{The committed tools}\label{t2ann:sec:outils}

Every \mes{} or \calc{} number of this volume comes from a committed tool. \Cref{t2ann:tab:outils} lists them, sorted by lot, each with the first line of its own description, quoted as written (in French), and its number of lines. The table is generated by \fichier{memoire/outils/tome2_annexes.py} from \texttt{git ls-files}: it follows the repository.

\begingroup
\captionof{table}[The tools of the engine and of its measurements]{The tools of the engine and of its measurements, by lot.}\label{t2ann:tab:outils}
\input{\Racine/tome2/annexes/outils-moteur}
\endgroup

\section{The figures and their data}\label{t2ann:sec:figures}

Every pgfplots graph of the volume reads a committed file; every TikZ diagram is a file under \fichier{memoire/figures/tome2/}. \Cref{t2ann:tab:figures} gives, for every figure command that reads data, the files it reads; the tool that generates it exits with an error if one of them is missing.

\begingroup
\captionof{table}[The figures and the files they read]{The figures of Volume~II and the data files they read.}\label{t2ann:tab:figures}
\input{\Racine/tome2/annexes/figures-tome2}
\endgroup

\section{The logs}\label{t2ann:sec:journaux}

The raw logs, CSV files and notes of every lot are committed under \fichier{mesures/moteur/}. \Cref{t2ann:tab:journaux} gives their extent, folder by folder (the tools, counted in \cref{t2ann:tab:outils}, are excluded). Heavy files stay outside the repository: the \texttt{.trc} files and their broken copies under \fichier{~/dev/modeles/trc/}, the outputs of the passes (\texttt{.kld}, \texttt{.npz}) under \fichier{~/dev/essais/troncons/}, the built programs under \fichier{~/dev/cibles/}; the logs keep their SHA-256 fingerprints and the commands.

\begingroup
\captionof{table}[The committed logs of the engine]{The committed logs of the engine, by folder.}\label{t2ann:tab:journaux}
\input{\Racine/tome2/annexes/journaux-moteur}
\endgroup

\begin{faux}
\begin{itemize}
\item \textbf{The first line of a description is not the tool}: it says what its author meant it to do. The chapters cite tools by what they produced, with their logs.
\item \textbf{The witness of the options compares names, not meanings}: a default changed in the code without \cref{t2ann:tab:options} changing would pass. The defaults of the table were reread in \fichier{moteur/crates/trc-ordo/src/lib.rs} (\texttt{Config::nouvelle}) and in \texttt{main.rs} on 16 October 2026.
\item \textbf{The heavy files outside the repository} are reproducible only through their fingerprints and the commands in the logs; their folder may be freed by Stone.
\end{itemize}
\end{faux}
